\documentclass[aspectratio=169,11pt]{beamer}
\usepackage[T1]{fontenc}
\usepackage{lmodern}
\usepackage[sfdefault,lining,tabular]{FiraSans}
\usepackage[scale=0.85]{FiraMono}
\usepackage{microtype}
\usepackage{array}
\usepackage{booktabs}
\usepackage{tabularx}
\usepackage{colortbl}
\usepackage{listings}
\usepackage{tikz}
\usetikzlibrary{arrows.meta,positioning,calc,shapes.geometric}

\usetheme[sectionpage=none,progressbar=frametitle,numbering=none]{metropolis}

% Color palette
\definecolor{Ink}{HTML}{14263C}
\definecolor{Muted}{HTML}{52657A}
\definecolor{Teal}{HTML}{00685E}
\definecolor{Pale}{HTML}{EAF5F2}
\definecolor{Rust}{HTML}{A2382C}
\definecolor{RustBg}{HTML}{FDF1F0}
\definecolor{CodeBg}{HTML}{F4F6F9}
\definecolor{Border}{HTML}{D6DFE8}
\definecolor{RowTint}{HTML}{F7F9FB}

\setbeamercolor{normal text}{fg=Ink,bg=white}
\setbeamercolor{frametitle}{fg=white,bg=Ink}
\setbeamercolor{structure}{fg=Teal}
\setbeamercolor{progress bar}{fg=Teal,bg=Border}
\setbeamercolor{footline}{fg=Muted,bg=white}

\setbeamerfont{title}{size=\Huge,series=\bfseries}
\setbeamerfont{frametitle}{size=\large,series=\bfseries}
\setbeamersize{text margin left=1cm,text margin right=1cm}

\setbeamertemplate{itemize item}{\color{Teal}\rule[0.2ex]{0.65ex}{0.65ex}}
\setlength{\leftmargini}{1.1em}

\makeatletter
\setlength{\metropolis@progressinheadfoot@linewidth}{1pt}
\makeatother

\setbeamertemplate{footline}{%
  \begin{beamercolorbox}[wd=\paperwidth,leftskip=1cm,rightskip=1cm]{footline}%
    \color{Border}\rule{\dimexpr\paperwidth-2cm\relax}{0.4pt}\par\vspace{2pt}%
    {\scriptsize\color{Muted}\secname}%
    \hfill%
    {\scriptsize\color{Muted}\insertframenumber\,/\,\inserttotalframenumber}%
    \vspace{4pt}%
  \end{beamercolorbox}%
}

\setbeameroption{hide notes}

% Code listings style
\lstset{
  language=C,
  basicstyle=\fontsize{6.2}{7.8}\selectfont\ttfamily\color{Ink},
  keywordstyle=\color{Teal}\bfseries,
  commentstyle=\color{Muted}\itshape,
  stringstyle=\color{Rust},
  backgroundcolor=\color{CodeBg},
  frame=leftline,
  framerule=2.5pt,
  rulecolor=\color{Ink!40},
  framesep=5pt,
  xleftmargin=7.5pt,
  aboveskip=1pt,
  belowskip=2pt,
  columns=fullflexible,
  keepspaces=true,
  showstringspaces=false,
  escapeinside={(*@}{@*)}
}

% Diagram styles
\tikzset{
  stage/.style={
    draw=Teal,
    fill=Pale,
    line width=0.6pt,
    rounded corners=2pt,
    align=center,
    minimum height=8.5mm,
    inner xsep=5pt,
    inner ysep=3pt,
    font=\footnotesize\color{Ink}
  },
  appstage/.style={
    stage,
    draw=Ink!50,
    fill=RowTint
  },
  rejected/.style={
    stage,
    draw=Rust,
    fill=RustBg
  },
  transition/.style={
    -{Stealth[length=1.8mm]},
    semithick,
    draw=Ink
  }
}

% Reusable helpers
\newcommand{\code}[1]{\texttt{#1}}
\newcommand{\bugline}[1]{\textcolor{Rust}{\bfseries #1}}
\newcommand{\extlink}[2]{\href{#1}{\textcolor{Teal}{#2}}}
\newcommand{\pinnedcode}[2]{\extlink{https://github.com/espressif/esp-idf/blob/7ad199d14a5d0fd2f61193c0f8a51a5d2bfa2c00/components/esp_https_ota/src/esp_https_ota.c#1}{#2}}
\newcommand{\pinnedhdr}[2]{\extlink{https://github.com/espressif/esp-idf/blob/7ad199d14a5d0fd2f61193c0f8a51a5d2bfa2c00/components/esp_https_ota/include/esp_https_ota.h#1}{#2}}
\newcommand{\examplecode}[2]{\extlink{https://github.com/espressif/esp-idf/blob/07bfc09d0c8837b4d149e1151c57635a0c02755a/examples/system/ota/simple_ota_example/main/simple_ota_example.c#1}{#2}}

\newcommand{\passpill}[1]{\colorbox{Pale}{\textcolor{Teal}{\bfseries\footnotesize\,#1\,}}}
\newcommand{\failpill}[1]{\colorbox{RustBg}{\textcolor{Rust}{\bfseries\footnotesize\,#1\,}}}

\newcommand{\srcbar}[2]{%
  \noindent{\footnotesize\ttfamily\color{Ink}#1}\hfill{\scriptsize\color{Muted}#2}\par\vspace{2pt}%
}

\newcommand{\accentcard}[4]{%
  \par\noindent\begin{tikzpicture}[baseline=(card.north)]
    \node[
      fill=#2,
      rounded corners=2pt,
      inner xsep=9pt,
      inner ysep=5pt,
      text width=\dimexpr\linewidth-18pt\relax,
      align=flush left,
      path picture={
        \fill[#1] (path picture bounding box.south west) rectangle
                  ([xshift=3pt]path picture bounding box.north west);
      }
    ] (card) {%
      {\small\bfseries\color{#1}#3}\par\vspace{2pt}%
      {\footnotesize\color{Ink}\setlength{\parskip}{2pt}#4}%
    };
  \end{tikzpicture}\par%
}

\newcommand{\pillarcard}[3]{%
  \noindent\begin{tikzpicture}[baseline=(card.north)]
    \node[
      fill=RowTint,
      draw=Border,
      line width=0.5pt,
      rounded corners=2pt,
      inner xsep=9pt,
      inner ysep=8pt,
      minimum height=4.05cm,
      text width=\dimexpr\linewidth-18.5pt\relax,
      align=flush left,
      path picture={
        \fill[Teal] (path picture bounding box.north west) rectangle
                    ([yshift=-3pt]path picture bounding box.north east);
      }
    ] (card) {%
      \parbox[t][3.45cm][t]{\linewidth}{%
        \rightskip=0pt plus 1fil\hyphenpenalty=10000\exhyphenpenalty=10000
        {\Large\bfseries\color{Teal}#1}\par\vspace{4pt}%
        {\normalsize\bfseries\color{Ink}#2}\par\vspace{5pt}%
        {\footnotesize\color{Ink}#3}%
      }%
    };
  \end{tikzpicture}%
}

\newcommand{\agendabox}[2]{%
  \begin{tikzpicture}[baseline=(box.north)]
    \node[
      fill=RowTint,
      draw=Border,
      line width=0.4pt,
      rounded corners=2pt,
      inner xsep=7pt,
      inner ysep=5.5pt,
      minimum height=12mm,
      text width=\dimexpr\linewidth-14.4pt\relax,
      align=flush left
    ] (box) {%
      {\scriptsize\bfseries\color{Teal}#1}\par\vspace{1pt}%
      {\footnotesize\color{Ink}#2}%
    };
  \end{tikzpicture}%
}

\title{ESP-IDF HTTPS OTA}
\subtitle{Failure handling, caller control, and charger recovery}
\author{Zaptec embedded software case}
\date{Assigned SDK revision 7ad199d14a}

\begin{document}

% Slide 1: Cover
\begin{frame}[plain]
  \begin{tikzpicture}[remember picture,overlay]
    \fill[Ink] (current page.north west) rectangle
               ([xshift=5mm,yshift=-0.65\paperheight]current page.north west);
    \fill[Teal] ([xshift=0mm,yshift=-0.65\paperheight]current page.north west) rectangle
                ([xshift=5mm]current page.south west);
  \end{tikzpicture}
  \vspace*{4mm}
  {\normalsize\bfseries\color{Teal}Zaptec embedded software case}\par\vspace{3mm}
  {\Huge\bfseries\color{Ink}ESP-IDF HTTPS OTA}\par\vspace{2mm}
  {\large\color{Muted}Failure handling, caller control, and charger recovery}\par\vspace{4mm}
  {\color{Teal}\rule{28mm}{1.5pt}}\par\vspace{4mm}
  {\small\color{Ink}
    Assigned SDK revision \pinnedcode{}{7ad199d14a}\hspace{8mm}
    Caller example revision \examplecode{}{07bfc09d0c}
  }\par\vfill
  \noindent
  \begin{minipage}[t]{0.235\textwidth}\agendabox{01}{Failure modes}\end{minipage}\hfill
  \begin{minipage}[t]{0.235\textwidth}\agendabox{02}{QEMU results}\end{minipage}\hfill
  \begin{minipage}[t]{0.235\textwidth}\agendabox{03}{Caller control}\end{minipage}\hfill
  \begin{minipage}[t]{0.235\textwidth}\agendabox{04}{Hardware checks}\end{minipage}
  \vspace*{2mm}
  \note{Target time is 45 seconds. Start with what the charger owner sees. Success means the new firmware boots and passes its health checks. A failed update must leave the charger safe and reachable.}
\end{frame}

\section{What can go wrong?}

% Slide 2: Overview flow
\begin{frame}{An OTA update has five stages}
  \begin{center}
    \begin{tikzpicture}[node distance=5mm]
      \node[stage] (download) {Download\\and write};
      \node[stage,right=of download] (validate) {Check\\image};
      \node[stage,right=of validate] (select) {Select\\boot slot};
      \node[appstage,right=of select] (reboot) {Reboot};
      \node[appstage,right=of reboot] (confirm) {Confirm\\health};
      \draw[transition] (download) -- (validate);
      \draw[transition] (validate) -- (select);
      \draw[transition] (select) -- (reboot);
      \draw[transition] (reboot) -- (confirm);
      \draw[Teal,semithick] (download.south west) -- ++(0,-2mm) -| (select.south east);
      \node[below=2.5mm,font=\footnotesize] at ($(download.south west)!0.5!(select.south east)$)
        {\pinnedcode{\#L381}{\code{esp\_https\_ota()}}};
      \draw[Ink!60,semithick] (reboot.south west) -- ++(0,-2mm) -| (confirm.south east);
      \node[below=2.5mm,font=\footnotesize,color=Ink] at ($(reboot.south west)!0.5!(confirm.south east)$)
        {Charger firmware};
    \end{tikzpicture}
  \end{center}
  \vspace{3mm}
  \noindent
  \begin{minipage}[t]{0.485\textwidth}
    \accentcard{Teal}{Pale}{What the SDK does}{%
      The wrapper \pinnedcode{\#L381}{\code{esp\_https\_ota()}} downloads the image, checks it, and selects the next boot slot.%
    }
  \end{minipage}\hfill
  \begin{minipage}[t]{0.485\textwidth}
    \accentcard{Ink}{RowTint}{What the charger does}{%
      The example \examplecode{\#L115}{\code{simple\_ota\_example.c}} reboots as soon as \code{ESP\_OK} returns. The charger must decide when a reboot is safe and run first-boot checks.%
    }
  \end{minipage}
  \note{Target time is 60 seconds. Show what the API covers and what the product owns. The example restarts right away when the wrapper returns success. Safe charging pauses and saved data safety belong to the charger firmware.}
\end{frame}

% Slide 3: Bug 1
\begin{frame}[fragile]{Boot selection fails, but OTA reports success}
  \begin{columns}[T,onlytextwidth]
    \begin{column}{0.53\textwidth}
      \srcbar{\code{esp\_https\_ota\_finish()}}{\pinnedcode{\#L340-L366}{lines 340 to 366}}
\begin{lstlisting}
esp_err_t err = ESP_OK;
// Line 344 sets outer err inside switch.
if ((err == ESP_OK) &&
    (handle->state == ESP_HTTPS_OTA_SUCCESS)) {
    (*@\bugline{esp\_err\_t err = esp\_ota\_set\_boot\_partition(}@*)
        (*@\bugline{handle->update\_partition);}@*)
    // Line 360 hides outer err; outer stays 0.
}
return err;
\end{lstlisting}
    \end{column}
    \begin{column}{0.44\textwidth}
      \accentcard{Rust}{RustBg}{Injected boot failure}{%
        Boot selection returned \code{ESP\_ERR\_FLASH\_OP\_FAIL} (\code{24577}). The OTA call returned \code{ESP\_OK} (\code{0}).%
      }
      \vspace{2mm}
      \accentcard{Ink}{RowTint}{Why this happens}{%
        An inner \code{err} hides the outer one at \pinnedcode{\#L360}{line 360}. The function logs the error and returns \code{ESP\_OK}.\par
        Issue \extlink{https://github.com/espressif/esp-idf/issues/12635}{\#12635} reports this bug. The \extlink{https://github.com/espressif/esp-idf/commit/d778b07b2213370e1603d1aadf7b60393b3f7569}{v5.1.3} fix updates the outer variable.%
      }
    \end{column}
  \end{columns}
  \note{Target time is 2 minutes. Point out the outer and inner err variables in esp\_https\_ota\_finish. The fix writes to the outer variable. Notice esp\_ota\_set\_boot\_partition validates the image a second time, so a bootloader\_mmap failure under memory or MMU pressure (issue 16968) can trigger this bug in practice. The injected test shows false success. Physical flash failure and next-boot behavior on real hardware are still untested. Upstream issue 12635 reports the same bug.}
\end{frame}

% Slide 4: Bug 2
\begin{frame}[fragile]{Setup loses the exact error}
  \begin{columns}[T,onlytextwidth]
    \begin{column}{0.55\textwidth}
      \srcbar{\code{esp\_https\_ota()}}{\pinnedcode{\#L392-L396}{lines 392 to 396}}
\begin{lstlisting}
esp_https_ota_handle_t https_ota_handle = NULL;
esp_err_t err = esp_https_ota_begin(
    &ota_config, &https_ota_handle);
if (https_ota_handle == NULL) {
    (*@\bugline{return ESP\_FAIL;}@*) // Drops err at line 395
}
\end{lstlisting}
    \end{column}
    \begin{column}{0.42\textwidth}
      \accentcard{Rust}{RustBg}{Injected memory failure}{%
        The test expected \code{ESP\_ERR\_NO\_MEM} (\code{257}). The wrapper returned \code{ESP\_FAIL} (\code{-1}).%
      }
      \vspace{2.5mm}
      \accentcard{Ink}{RowTint}{Why this matters}{%
        The header \pinnedhdr{\#L53-L60}{\code{esp\_https\_ota.h}} (lines 53 to 60) promises the exact error code.\par
        Retry code cannot tell memory or argument errors from a general failure.%
      }
    \end{column}
  \end{columns}
  \note{Target time is 75 seconds. Look at esp\_https\_ota right after begin fails. Returning the exact error helps the caller pick a smart retry rule. The test forces calloc inside the component to fail.}
\end{frame}

% Slide 5: Bug 3
\begin{frame}[fragile]{HTTP errors start a flash update}
  \begin{center}
    \begin{tikzpicture}[node distance=5mm]
      \node[rejected] (http) {HTTP 404\\or 500};
      \node[appstage,right=of http] (accepted) {Connection\\accepted};
      \node[stage,right=of accepted] (flash) {\code{esp\_ota\_begin()}\\erases slot};
      \node[rejected,right=of flash] (validate) {Image check fails\\(\code{5379})};
      \draw[transition] (http) -- (accepted);
      \draw[transition] (accepted) -- (flash);
      \draw[transition] (flash) -- (validate);
    \end{tikzpicture}
  \end{center}
  \vspace{3mm}
  \begin{columns}[T,onlytextwidth]
    \begin{column}{0.56\textwidth}
      \srcbar{\code{\_http\_handle\_response\_code()}}{\pinnedcode{\#L60-L88}{lines 60 to 88, 271}}
\begin{lstlisting}
// Redirect and 401 handling omitted here.
if (process_again(status_code)) {
    // Comment says redirects only (line 74),
    // yet process_again includes 401 (line 52).
}
(*@\bugline{return ESP\_OK;}@*) // Reached for HTTP 404 and 500
\end{lstlisting}
    \end{column}
    \begin{column}{0.41\textwidth}
      \accentcard{Rust}{RustBg}{Flash wear and endless retries}{%
        HTTP 404 and 500 must fail before any flash erase, with zero calls to \code{esp\_ota\_begin()}.\par
        HTTP 401 also retries with no limit inside \pinnedcode{\#L109}{\code{\_http\_connect()}}.%
      }
    \end{column}
  \end{columns}
  \note{Target time is 105 seconds. Show \_http\_handle\_response\_code and the first branch of esp\_https\_ota\_perform. Note that esp\_ota\_begin with OTA\_SIZE\_UNKNOWN erases the passive slot before checking the HTTP body. Also note that \_http\_connect retries HTTP 401 with no limit and does not enforce redirect limits in either 7ad199d14a or v6.1.1. Checking only the final error would hide the flash erase.}
\end{frame}

% Slide 6: Bug 4
\begin{frame}[fragile]{The wrapper has no total time limit}
  \begin{columns}[T,onlytextwidth]
    \begin{column}{0.54\textwidth}
      \srcbar{\code{esp\_https\_ota\_perform()}}{\pinnedcode{\#L285-L306}{lines 285 to 306, 398 to 403}}
\begin{lstlisting}[basicstyle=\fontsize{6.1}{7.6}\selectfont\ttfamily\color{Ink}]
data_read = esp_http_client_read(
    handle->http_client, ...);
if (data_read == 0) {
    bool done =
        esp_https_ota_is_complete_data_received(
            https_ota_handle);
    if ((errno == ENOTCONN || errno == ECONNRESET ||
         errno == ECONNABORTED) && !done) {
        return ESP_FAIL;
    }
    if (!done) {
        (*@\bugline{return ESP\_ERR\_HTTPS\_OTA\_IN\_PROGRESS;}@*)
    }
}
\end{lstlisting}
    \end{column}
    \begin{column}{0.43\textwidth}
      \accentcard{Rust}{RustBg}{Loop with no limit}{%
        A read timeout leaves \code{errno} at zero in \code{7ad199d14a} and returns \code{-ESP\_ERR\_HTTP\_EAGAIN} in \code{v6.1.1}.\par
        \pinnedcode{\#L398-L403}{\code{esp\_https\_ota()}} loops in \code{while (1)} with no deadline.%
      }
      \vspace{2.5mm}
      \accentcard{Ink}{RowTint}{Stall test and issue \#6364}{%
        Issue \extlink{https://github.com/espressif/esp-idf/issues/6364}{\#6364} reports a 10-hour stall. In QEMU, \code{7ad199d14a} and \code{v6.1.1} spin past the limit; the local fix returns \code{263}.%
      }
    \end{column}
  \end{columns}
  \note{Target time is 75 seconds. Walk through esp\_https\_ota\_perform at lines 285 to 306 and the while loop in esp\_https\_ota at lines 398 to 403. In QEMU, a stalled 200 stream triggers 14 read timeouts on 7ad199d14a and 15 on v6.1.1 until the test server closes after 8 seconds. The local fix stops with ESP\_ERR\_TIMEOUT (263) after 4 timeouts.}
\end{frame}

% Slide 7: QEMU results
\begin{frame}{One baseline and five bug tests}
  \begingroup
  \setlength{\aboverulesep}{0pt}
  \setlength{\belowrulesep}{0pt}
  \renewcommand{\arraystretch}{1.05}
  \footnotesize
  \begin{tabularx}{\textwidth}{ X c c c }
    \toprule
    \rowcolor{Ink!10}
    \textbf{Case} & \textbf{Assigned SDK} & \textbf{Local fixes} & \textbf{v6.1.1} \\
    \midrule
    Good download and image check & \passpill{Pass} (\code{0}) & \passpill{Pass} & \passpill{Pass} \\
    \rowcolor{RowTint}
    Boot selection fails (injected) & \failpill{Fail} (\code{0}, expected \code{24577}) & \passpill{Pass} & \passpill{Pass} \\
    Memory allocation fails (injected) & \failpill{Fail} (\code{-1}, expected \code{257}) & \passpill{Pass} & \passpill{Pass} \\
    \rowcolor{RowTint}
    HTTP 404 before OTA starts & \failpill{Fail} (\code{5379}, \code{flash=1}) & \passpill{Pass} & \passpill{Pass} \\
    HTTP 500 before OTA starts & \failpill{Fail} (\code{5379}, \code{flash=1}) & \passpill{Pass} & \passpill{Pass} \\
    \rowcolor{RowTint}
    Stalled download (no deadline) & \failpill{Fail} (\code{-1}, expected \code{263}) & \passpill{Pass} & \failpill{Fail} (\code{-1}) \\
    \bottomrule
  \end{tabularx}
  \endgroup
  \par\vspace{3mm}
  \noindent
  \begin{minipage}[t]{0.485\textwidth}
    \accentcard{Teal}{Pale}{How I checked}{%
      A language model flagged bug paths. I checked each path in \code{7ad199d14a} plus \code{flash\_begin}, HTTP status, and timeouts.%
    }
  \end{minipage}\hfill
  \begin{minipage}[t]{0.485\textwidth}
    \accentcard{Ink}{RowTint}{Test limits}{%
      Code at \code{7ad199d14a} supports ESP32 and ESP32-S2. QEMU runs one core over HTTP with disabled watchdogs.%
    }
  \end{minipage}
  \note{Target time is 90 seconds. Highlight the QEMU proof. Code 5379 (ESP\_ERR\_OTA\_VALIDATE\_FAILED) with flash\_begin equal to 1 proves esp\_ota\_begin ran on HTTP 404 and 500. Stalled download fails on both 7ad199d14a and v6.1.1 because esp\_https\_ota loops until the server closes the socket. Also note the single-core, HTTP, watchdog, and chip limits.}
\end{frame}

% Slide 8: Bug summary table
\begin{frame}{Bugs, fixes, and upstream reports}
  \begingroup
  \setlength{\aboverulesep}{0pt}
  \setlength{\belowrulesep}{0pt}
  \renewcommand{\arraystretch}{1.12}
  \footnotesize
  \begin{tabularx}{\textwidth}{
    >{\raggedright\arraybackslash}p{0.31\textwidth}
    >{\raggedright\arraybackslash}p{0.22\textwidth}
    >{\raggedright\arraybackslash}p{0.18\textwidth}
    >{\raggedright\arraybackslash}X }
    \toprule
    \rowcolor{Ink!10}
    \textbf{Bug} & \textbf{Pinned code} & \textbf{Fix verified} & \textbf{ESP-IDF issue} \\
    \midrule
    False success on boot selection
      & \pinnedcode{\#L360}{Line 360}
      & \extlink{https://github.com/espressif/esp-idf/commit/d778b07b2213370e1603d1aadf7b60393b3f7569}{v5.1.3} (5.1 branch)
      & \extlink{https://github.com/espressif/esp-idf/issues/12635}{\#12635} (exact) \\
    \rowcolor{RowTint}
    Setup error lost
      & \pinnedcode{\#L393-L395}{Lines 393 to 395}
      & \extlink{https://github.com/espressif/esp-idf/blob/v6.1.1/components/esp_https_ota/src/esp_https_ota.c\#L1099-L1102}{v6.1.1}
      & No exact report \\
    Flash erase on HTTP 404 or 500
      & \pinnedcode{\#L60-L110}{Lines 60 to 110}, \pinnedcode{\#L271}{271}
      & \extlink{https://github.com/espressif/esp-idf/commit/3c7991813d8a8f47009305a63d7d8d0d422fb8fb}{v4.4} source
      & \extlink{https://github.com/espressif/esp-idf/issues/7058}{\#7058}, \extlink{https://github.com/espressif/esp-idf/issues/8195}{\#8195} (related) \\
    \rowcolor{RowTint}
    No total time limit for OTA
      & \pinnedcode{\#L285-L306}{Lines 285 to 306}, \pinnedcode{\#L398-L403}{398 to 403}
      & Unfixed in \code{v6.1.1}; local fix
      & \extlink{https://github.com/espressif/esp-idf/issues/6364}{\#6364} \\
    \bottomrule
  \end{tabularx}
  \endgroup
  \vspace{2.5mm}
  \accentcard{Teal}{Pale}{What passed}{%
    Bugs 1 to 3 passed locally and on \code{v6.1.1}. Bug 4 still fails on \code{v6.1.1} and passes with the local no-progress and total time limit.%
  }
  \note{Target time is 45 seconds. The links point to esp\_https\_ota.c at the assigned commit. The third column lists versions with tested fixes. Bugs 1 to 3 passed on v6.1.1, while Bug 4 (stalled download with no deadline) still fails on v6.1.1 and passes only with our local fix.}
\end{frame}

\section{How should the application call it?}

% Slide 9: Example analysis
\begin{frame}[fragile]{Why the assigned example is not enough for a charger}
  \begin{columns}[T,onlytextwidth]
    \begin{column}{0.52\textwidth}
      \srcbar{\code{simple\_ota\_example.c}}{\examplecode{\#L114-L123}{lines 114 to 123, 156 to 164}}
\begin{lstlisting}[basicstyle=\fontsize{6.2}{7.8}\selectfont\ttfamily\color{Ink}]
esp_err_t ret = esp_https_ota(&config);
if (ret == ESP_OK) {
    (*@\bugline{esp\_restart();}@*) // Reboots during charging
} else {
    ESP_LOGE(TAG, "Firmware upgrade failed");
}
(*@\bugline{while (1) \{}@*) // Stuck task; no retry
    (*@\bugline{vTaskDelay(1000 / portTICK\_PERIOD\_MS);}@*)
(*@\bugline{\}}@*)
\end{lstlisting}
    \end{column}
    \begin{column}{0.45\textwidth}
      \accentcard{Ink}{RowTint}{Problems in the example}{%
        The example reboots right away on success, gets stuck on failure, and never confirms the new image.\par
        Startup (\examplecode{\#L157-L163}{lines 157 to 163}) erases NVS on a version mismatch, which can wipe credentials after a rollback.%
      }
      \vspace{2.5mm}
      \accentcard{Rust}{RustBg}{Can the wrapper be safe enough?}{%
        No. The one-shot wrapper hides boot and setup errors, skips checks before flash erase, and has no total time limit.%
      }
    \end{column}
  \end{columns}
  \note{Target time is 75 seconds. Answer Questions 2 and 3 from the brief. The wrapper esp\_https\_ota cannot be safe enough as written. Also, simple\_ota\_example.c calls esp\_restart during active charging, enters an idle loop on failure, and erases NVS on a version mismatch.}
\end{frame}

% Slide 10: Staged API sketch
\begin{frame}[fragile]{Step-by-step calls give more control}
  \begin{columns}[T,onlytextwidth]
    \begin{column}{0.55\textwidth}
      \srcbar{Step-by-step caller sketch}{\pinnedcode{\#L210}{lines 210, 293}}
\begin{lstlisting}[basicstyle=\fontsize{6.1}{7.6}\selectfont\ttfamily\color{Ink}]
err = esp_https_ota_begin(&cfg, &handle);
if (err != ESP_OK) return err;
err = esp_https_ota_get_img_desc(handle, &info);
if (err == ESP_OK && !valid_header(&info))
    err = ESP_ERR_INVALID_VERSION;
if (err == ESP_OK) do {
    err = timed_out() ? ESP_ERR_TIMEOUT
        : esp_https_ota_perform(handle);
} while (err == ESP_ERR_HTTPS_OTA_IN_PROGRESS);
// Use abort() on error in v4.4 and newer.
esp_err_t finish_err = esp_https_ota_finish(handle);
return err == ESP_OK ? finish_err : err;
\end{lstlisting}
    \end{column}
    \begin{column}{0.42\textwidth}
      \accentcard{Teal}{Pale}{Check before erase}{%
        Call \code{esp\_https\_ota\_get\_img\_desc()} before any flash erase.%
      }
      \vspace{2mm}
      \accentcard{Ink}{RowTint}{Time-limited loop}{%
        Run in one worker task. Stop \code{esp\_https\_ota\_perform()} when time runs out.%
      }
      \vspace{2mm}
      \accentcard{Ink}{RowTint}{Safe reboot}{%
        Reboot only after finish succeeds and charging is safe.%
      }
    \end{column}
  \end{columns}
  \note{Target time is 105 seconds. Call esp\_https\_ota\_begin, then esp\_https\_ota\_get\_img\_desc before esp\_ota\_begin erases flash, then esp\_https\_ota\_perform in a timed loop. Perform reaches success only when esp\_https\_ota\_is\_complete\_data\_received returns true at line 293. On 7ad199d14a, the application must also serialize OTA with a lock (issue 18127) and keep the OTA task stack in internal RAM (issue 10437). Finish is the only cleanup call on 7ad199d14a, and v4.4 and newer add esp\_https\_ota\_abort.}
\end{frame}

% Slide 11: Staged API limitations
\begin{frame}{The caller still needs a fixed SDK}
  \begin{columns}[T,onlytextwidth]
    \begin{column}{0.50\textwidth}
      \accentcard{Ink}{RowTint}{Limits inside the pinned SDK}{%
        The step-by-step API uses the same HTTP and boot-selection code as \code{esp\_https\_ota()}.\par
        Code at \code{7ad199d14a} also allows HTTPS downgrades (\extlink{https://github.com/espressif/esp-idf/commit/eef3b0bf}{eef3b0bf}) and blocks certificate bundles (\extlink{https://github.com/espressif/esp-idf/issues/5504}{\#5504}).\par
        \code{begin()} and \pinnedcode{\#L231-L245}{\code{get\_img\_desc()}} can block in read loops, where caller timers cannot stop them.%
      }
    \end{column}
    \begin{column}{0.47\textwidth}
      \accentcard{Teal}{Pale}{What must change}{%
        Use a fixed SDK or backport the fixes. Limit network waits and retries, then test the total time limit.\par
        Even newer SDKs need caller guards for redirect limits and open \code{esp\_ota\_ops} bugs.%
      }
      \vspace{2.5mm}
      \accentcard{Ink}{RowTint}{Retry rules}{%
        Retry short network errors with backoff. Stop right away on bad certificates or bad images.%
      }
    \end{column}
  \end{columns}
  \note{Target time is 90 seconds. One worker task keeps the charger responsive during downloads. Beyond the four core bugs, 7ad199d14a allows HTTPS downgrades (commit eef3b0bf) and blocks certificate bundles (issue 5504). Newer SDKs such as v6.1.1 still have unguarded redirect loops in \_http\_connect and open esp\_ota\_ops bugs (issues 19413 and 19439).}
\end{frame}

\section{How do we test charger OTA?}

% Slide 12: Hardware verification matrix
\begin{frame}{Test recovery on real hardware}
  \begingroup
  \setlength{\aboverulesep}{0pt}
  \setlength{\belowrulesep}{0pt}
  \renewcommand{\arraystretch}{1.12}
  \footnotesize
  \begin{tabularx}{\textwidth}{ >{\raggedright\arraybackslash}p{0.39\textwidth} >{\raggedright\arraybackslash}X }
    \toprule
    \rowcolor{Ink!10}
    \textbf{Fault case} & \textbf{Expected result} \\
    \midrule
    Network stalls, drops, and cut-off downloads
      & The update stops within the time limit. The old firmware still works. \\
    \rowcolor{RowTint}
    Bad certificate or unapproved image
      & The update is rejected. Current firmware stays usable. \\
    Power loss during flash write, boot selection, or first boot
      & A working image boots, the charger stays safe, and saved settings survive. \\
    \rowcolor{RowTint}
    Failed first-boot health check
      & The charger rolls back, and old firmware keeps its data. \\
    \bottomrule
  \end{tabularx}
  \endgroup
  \vspace{2.5mm}
  \accentcard{Teal}{Pale}{Before testing}{%
    Use production TLS, watchdog, and bootloader settings; test real heap load (\extlink{https://github.com/espressif/esp-idf/issues/16968}{\#16968}) and older field builds (\extlink{https://github.com/espressif/esp-idf/issues/9987}{\#9987}, \extlink{https://github.com/espressif/esp-idf/issues/17595}{\#17595}).%
  }
  \note{Target time is 105 seconds. Describe a hardware bench with network faults, power cuts, and charging checks. Include dual-core flash watchdog checks (issues 10890 and 14450), realistic firmware sizes under heap and MMU load (issue 16968), and upgrades from older field firmware (issues 9987 and 17595). Agree on update time limits, charging pause rules, and recovery goals before testing. QEMU disables watchdogs and runs one core over HTTP.}
\end{frame}

% Slide 13: Rollback flow
\begin{frame}{A failed first boot must roll back}
  \begin{center}
    \begin{tikzpicture}[node distance=9mm]
      \node[stage] (boot) {Boot new image};
      \node[stage,right=of boot] (health) {Wait for checks\\Run health tests};
      \coordinate[right=8mm of health.east] (branch);
      \node[stage,anchor=west] (confirm) at ([xshift=29mm,yshift=8.5mm]health.east)
        {Confirm\\new firmware};
      \node[rejected,anchor=west] (rollback) at ([xshift=29mm,yshift=-8.5mm]health.east)
        {Roll back to\\last good image};
      \draw[transition] (boot) -- (health);
      \draw[semithick,draw=Ink] (health.east) -- (branch);
      \draw[transition] (branch) |-
        node[pos=0.75,above,font=\footnotesize,color=Teal] {Pass} (confirm.west);
      \draw[transition] (branch) |-
        node[pos=0.75,below,font=\footnotesize,color=Rust] {Fail or reset} (rollback.west);
    \end{tikzpicture}
  \end{center}
  \vspace{2.5mm}
  \noindent
  \begin{minipage}[t]{0.485\textwidth}
    \accentcard{Teal}{Pale}{Confirm the new image}{%
      Turn on bootloader rollback. Call\par\vspace{1pt}
      {\fontsize{6.4}{8}\selectfont\code{esp\_ota\_mark\_app\_valid\_cancel\_rollback()}}\par\vspace{1pt}
      only after key charger checks pass.%
    }
  \end{minipage}\hfill
  \begin{minipage}[t]{0.485\textwidth}
    \accentcard{Rust}{RustBg}{Automatic rollback}{%
      If checks fail or the charger resets before confirmation, the bootloader boots the last good image next.%
    }
  \end{minipage}
  \note{Target time is 75 seconds. A healthy image calls esp\_ota\_mark\_app\_valid\_cancel\_rollback. A failed test requests rollback. Watch for open issue 19413, where a retry after a failed otadata write can return ESP\_OK without confirming the running firmware. A reset before confirmation also rolls back when a good old image exists. Test this on real hardware.}
\end{frame}

% Slide 14: Rollout policy
\begin{frame}{Test the reboot rule before a wide rollout}
  \begin{center}
    \begin{tikzpicture}[node distance=18mm]
      \node[stage] (s1) {\textbf{Stage 1}\\Lab tests};
      \node[stage,right=of s1] (s2) {\textbf{Stage 2}\\Small field group};
      \node[appstage,right=of s2] (s3) {\textbf{Stage 3}\\Full rollout};
      \draw[transition] (s1) -- node[above,font=\scriptsize,color=Teal] {Pass} (s2);
      \draw[transition] (s2) -- node[above,font=\scriptsize,color=Ink] {Healthy} (s3);
    \end{tikzpicture}
  \end{center}
  \vspace{3mm}
  \noindent
  \begin{minipage}[t]{0.485\textwidth}
    \accentcard{Teal}{Pale}{What to track}{%
      Track these items before a wider rollout:
      \begin{itemize}\setlength{\itemsep}{1pt}\setlength{\topsep}{2pt}\setlength{\parsep}{0pt}
        \item Reboot tests during idle and active charging
        \item Firmware version reported after reboot
        \item Time to confirmed boot and rollback rate
        \item Chargers that stop reporting
      \end{itemize}%
    }
  \end{minipage}\hfill
  \begin{minipage}[t]{0.485\textwidth}
    \accentcard{Ink}{RowTint}{When to pause}{%
      Agree on rollback, stall, and offline limits with product and operations teams before a wider rollout.\par
      Pause the rollout right away if the small group crosses any limit.%
    }
  \end{minipage}
  \note{Target time is 45 seconds. Agree on pause limits with product and operations teams. Always read the firmware version from the charger after it reboots.}
\end{frame}

% Slide 15: How I would ship OTA
\begin{frame}{How I would ship OTA}
  \vspace{1mm}
  \noindent
  \begin{minipage}[t]{0.315\textwidth}
    \pillarcard{01}{Fixed SDK}{%
      Pick an ESP-IDF version or patched component that passes all bug tests.%
    }
  \end{minipage}\hfill
  \begin{minipage}[t]{0.315\textwidth}
    \pillarcard{02}{Bounded caller}{%
      Run step-by-step calls in one task with time limits, few retries, and a safe reboot check.%
    }
  \end{minipage}\hfill
  \begin{minipage}[t]{0.315\textwidth}
    \pillarcard{03}{Hardware proof}{%
      Test safe charging, first-boot checks, and remote recovery on real hardware.%
    }
  \end{minipage}
  \vspace{3mm}
  \accentcard{Teal}{Pale}{Main goal}{%
    A failed update must always leave the charger safe, working, and reachable.%
  }
  \note{Target time is 90 seconds. Tie each point back to the four questions in the brief. Note that upgrading ESP-IDF fixes older bugs, yet newer releases still add regressions such as issues 13065 and 19439, so automated fault tests remain necessary. Invite questions on time limits, rollback checks, and fault tests.}
\end{frame}

\section{Sources and test run}

% Slide 16: Source and reproduction
\begin{frame}{Sources and test run}
  \begin{columns}[T,onlytextwidth]
    \begin{column}{0.50\textwidth}
      \accentcard{Ink}{RowTint}{Code and upstream links}{%
        Review sources and tracked issues:
        \begin{itemize}\setlength{\itemsep}{2pt}\setlength{\topsep}{2pt}\setlength{\parsep}{0pt}
          \item \pinnedcode{}{Assigned OTA source at \code{7ad199d14a}}
          \item \examplecode{}{Assigned caller example at \code{07bfc09d0c}}
          \item \extlink{https://github.com/espressif/esp-idf/issues/12635}{Boot-selection issue \#12635}
          \item \extlink{https://github.com/espressif/esp-idf/issues/7058}{HTTP status issues \#7058} and \extlink{https://github.com/espressif/esp-idf/issues/8195}{\#8195}
          \item \extlink{https://github.com/espressif/esp-idf/issues/6364}{Time limit issue \#6364}
        \end{itemize}%
      }
    \end{column}
    \begin{column}{0.47\textwidth}
      \accentcard{Teal}{Pale}{Run the QEMU tests}{%
        Run the test suite from the prepared repository:\par\vspace{4pt}
        \colorbox{white}{\parbox{\dimexpr\linewidth-6pt\relax}{\ttfamily\footnotesize\color{Ink}nix develop -c make test}}\par\vspace{5pt}
        Results are in \code{docs/qemu-results.md}.\par
        More issues are in \code{docs/related-issues.md}.%
      }
    \end{column}
  \end{columns}
  \note{Use this slide during questions and answers for commit hashes, issue numbers, or test commands.}
\end{frame}

\end{document}
